\documentclass[10pt,twocolumn]{article}
\usepackage[T1]{fontenc}
\usepackage[utf8]{inputenc}
\usepackage{lmodern}
\usepackage[margin=1.8cm,top=2.4cm,bottom=2cm,columnsep=0.6cm]{geometry}
\usepackage{fancyhdr}
\usepackage{titlesec}
\usepackage{booktabs}
\usepackage{amsmath,amssymb,amsfonts}
\usepackage{microtype}
\usepackage{xcolor}
\usepackage{mdframed}
\usepackage{parskip}
\usepackage{enumitem}
\usepackage{hyperref}

\definecolor{rulered}{RGB}{180,30,30}
\definecolor{boxgray}{RGB}{245,245,245}
\definecolor{keywordgray}{RGB}{100,100,100}

\pagestyle{fancy}
\fancyhf{}
\fancyhead[L]{\textsc{\small Claude Mag}}
\fancyhead[C]{\small\textit{Machine Learning Research Digest}}
\fancyhead[R]{\small 2026-07-01}
\fancyfoot[C]{\small\thepage}
\renewcommand{\headrulewidth}{0.8pt}

\titleformat{\section}{\normalsize\bfseries\color{rulered}}{}{0pt}{}%
  [\vspace{-4pt}\color{rulered}\rule{\columnwidth}{0.4pt}]
\titlespacing{\section}{0pt}{8pt}{4pt}

\hypersetup{colorlinks=true,urlcolor=rulered,linkcolor=black}

\begin{document}
\setlength{\parindent}{0pt}
\setlength{\parskip}{4pt}

{\small\color{keywordgray}\textbf{Keywords:}
  diffusion models \textbullet{}
  memorization \textbullet{}
  kernel density estimation \textbullet{}
  generalization theory}\par\vspace{4pt}

{\LARGE\bfseries\raggedright
  Local Coverage Governs Memorization in Diffusion Models}\par\vspace{4pt}

{\large\itshape\raggedright
  A Sharp Phase Transition Separates Memorized from Generalised Outputs---and
  It Depends Only on Neighbourhood Density}\par\vspace{6pt}

{\small\textbf{By} Claudia Merger \& Sebastian Goldt
  (SISSA / ICTP, Trieste)}\par\vspace{2pt}

{\small\color{keywordgray}arXiv:2606.14390 \textbullet{} Submitted 12 June 2026}
\par\vspace{2pt}\noindent{\color{rulered}\rule{\columnwidth}{0.6pt}}\vspace{6pt}

A new theoretical result explains why diffusion models memorise some training
examples and generalise on others---sometimes simultaneously. The criterion
is purely local: a training point is memorised if and only if its neighbourhood
in the data space contains too few other training samples for the model to
interpolate. In the high-dimensional limit, the transition is sharp and
its critical threshold scales as $n^{-1/k}$ with dataset size $n$ and
intrinsic data dimensionality $k$.

\begin{mdframed}[backgroundcolor=boxgray,linecolor=rulered,linewidth=1pt,
    innerleftmargin=6pt,innerrightmargin=6pt,
    innertopmargin=6pt,innerbottommargin=6pt]
\textbf{\small AT A GLANCE}\par\vspace{4pt}
\begin{description}[leftmargin=0pt,labelsep=4pt,itemsep=2pt,parsep=0pt,
    font=\small\bfseries]
  \item[Problem:] No principled theory predicts which training points a
    diffusion model will memorise versus generalise.
  \item[Why it matters:] Memorisation enables training-data extraction
    attacks, raises copyright concerns, and reveals a gap in our theoretical
    understanding of diffusion generative models.
  \item[Method:] Exploit the KDE interpretation of the diffusion score
    function to derive a closed-form local coverage criterion for
    memorisation, with a sharp phase transition in the high-$d$ limit.
  \item[Result:] The criterion correctly predicts memorised vs.\ novel
    samples on CIFAR-10 and CelebA without any model-specific information.
  \item[Evidence:] Validated on synthetic Gaussian mixtures ($d\in\{10,50,200\}$),
    CIFAR-10, and CelebA.
\end{description}
\end{mdframed}

\section{The Big Picture}

Diffusion models generate images by iteratively denoising a Gaussian noise
vector. When trained on a large dataset, most outputs are novel; but a small
fraction---particularly samples close to rare or isolated training images---are
effectively verbatim copies of training data. This memorisation is not a bug
in the model; it is a fundamental property of how the score function behaves
near sparse regions of the data distribution.

The community has documented memorisation empirically (cataloguing which prompts
or seeds trigger it) and proposed detection heuristics (membership inference,
attention-map analysis). What has been missing is a simple, principled answer
to the question: \emph{why does a particular training point get memorised?}

This paper answers that question. Leveraging the well-known connection between
score matching and kernel density estimation (KDE), the authors show that the
score function is dominated by the nearest training neighbours of any query
point. If a training point has too few neighbours within the relevant kernel
bandwidth, the score field around it points straight back to that point at
every noise level---ensuring the reverse diffusion ODE converges there.
In high dimensions, this transition is sharp.

\section{Why Existing Approaches Fall Short}

Empirical approaches study memorisation as a global property of the
trained model: ``this model memorises 3\% of its training set.'' They
cannot predict \emph{which} 3\% without running the model.

Attention-map analyses (Bright Ending Attention, BEA) detect memorisation
post hoc by examining the model's internal activations. They require
model access and do not provide a generalisable theory.

Membership inference attacks frame memorisation as a binary classification
problem, but they neither explain the mechanism nor provide a quantitative
criterion depending on data geometry.

\section{Core Mechanism}

The score function of a diffusion model at noise level $\sigma$ is
well-approximated (for overparameterised networks) by the KDE score:
\[
\hat s(x,\sigma) = \nabla_x \log \hat p_\sigma(x),
\quad
\hat p_\sigma(x) = \frac{1}{n}\sum_{i=1}^n \mathcal{N}(x;\,x_i,\,\sigma^2 I).
\]

Near a training point $x_j$, the score is:
\[
\hat s(x_j, \sigma)
\approx \frac{\sum_{i\neq j}(x_i - x_j)\,w_{ij}(\sigma)}
             {1 + \sum_{i\neq j} w_{ij}(\sigma)},
\quad
w_{ij}(\sigma) = e^{-\|x_i - x_j\|^2/2\sigma^2}.
\]
When $\sum_{i\neq j} w_{ij}(\sigma) \ll 1$ (no close neighbours), the score
at $x_j$ points approximately to zero---i.e., toward $x_j$ itself at all
scales. The reverse ODE then converges to $x_j$ from any initialisation
in its vicinity: memorisation.

\section{Technical Formulation}

\textbf{Memorisation criterion.}
Define the \emph{local coverage} of $x_j$ at scale $\sigma$:
\[
C_\sigma(x_j) = \frac{1}{n}\sum_{i\neq j}
  \exp\!\left(-\frac{\|x_i - x_j\|^2}{2\sigma^2}\right).
\]
Training point $x_j$ is \textbf{memorised} iff
$C_\sigma(x_j) \ll 1/n$ for all $\sigma \in [\sigma_\min, \sigma_\max]$
(the range of noise levels used during training).

\textbf{Phase transition in high dimensions.}
For data supported on a $k$-dimensional manifold of radius $R$, the
number of training samples within ball $B(x_j, \sigma)$ is
$\bar n(\sigma) \approx n\,(\sigma/R)^k$. Memorisation occurs iff
$\bar n(\sigma) = O(1)$, i.e., iff $\sigma < \sigma^*$ where:
\[
\sigma^* \sim R\,n^{-1/k}.
\]
This is the \textbf{critical radius}: points $x_j$ with no neighbours
within $\sigma^*$ are always memorised; points with many neighbours
within $\sigma^*$ always generalise. The transition is \emph{sharp} in
the limit $d \to\infty$, $n\to\infty$ with $k$ fixed.

\textbf{Predictor.}
The paper proposes a practical memorisation predictor: compute
$C_{\hat\sigma^*}(x_j)$ in the embedding space of a pre-trained feature
extractor and threshold at $C^* = 1/n$. This requires no model
access---only the dataset.

\section{Learning or Inference Procedure}

The analysis is purely theoretical; no new training procedure is introduced.
The practical predictor works as follows:
\begin{enumerate}[itemsep=2pt,parsep=0pt]
  \item Embed all training points with a feature extractor $\phi$.
  \item Estimate the intrinsic dimension $k$ of the training set via
    nearest-neighbour methods.
  \item Compute $\sigma^* = R\,n^{-1/k}$ where $R$ is the dataset diameter.
  \item For each training point $x_j$, compute $C_{\sigma^*}(x_j)$.
  \item Predict memorisation if $C_{\sigma^*}(x_j) < 1/n$.
\end{enumerate}

\section{What the Guarantee Says}

In the high-dimensional limit ($d\to\infty$ with $k/d\to0$), the
transition is shown to be sharp: for any $\varepsilon>0$, there exist
constants $c_1,c_2>0$ such that
\[
\Pr[\text{memorised} \mid C_\sigma(x_j) < c_1/n] \to 1,
\]
\[
\Pr[\text{generalised} \mid C_\sigma(x_j) > c_2/n] \to 1,
\]
as $n,d\to\infty$. Memorisation probability varies continuously between
these extremes at finite $n$.

\section{Experimental Findings}

\begin{itemize}[itemsep=2pt,parsep=0pt]
  \item \textbf{Synthetic:} On Gaussian mixtures with controlled
    inter-point distance, the empirical memorisation fraction matches
    the theoretical prediction to within measurement noise at
    $d\in\{10,50,200\}$.
  \item \textbf{CIFAR-10:} The local coverage criterion identifies
    memorised samples (confirmed by running inference from fixed seeds)
    with $>85\%$ precision and $>80\%$ recall without any model access.
  \item \textbf{CelebA:} Rare attribute combinations (e.g., red hair +
    glasses) are correctly predicted as memorised; common attributes
    (e.g., brown eyes) are correctly predicted as generalised.
\end{itemize}

\section{Ablations and Interpretation}

\textbf{Feature space:} Using raw pixel distances yields weaker
predictions than CLIP embeddings ($\Delta$-precision $\approx +15\%$),
confirming that perceptual distance is the relevant metric.

\textbf{Kernel bandwidth:} Results are robust to $\sigma$ within a
factor of 2 of $\sigma^*$, consistent with the sharp transition theory.

\textbf{Dataset size:} As $n$ increases, the fraction of memorised
samples decreases as $n^{-1/k}$, matching the theoretical scaling.

\section*{Reference}

Claudia Merger, Sebastian Goldt.
\textbf{Local Coverage Governs Memorization in Diffusion Models}.
arXiv:2606.14390, June 2026.
\url{https://arxiv.org/abs/2606.14390}

\end{document}
